\item Usted tiene un interés particular en los motores de automóviles, por lo que ha asegurado un puesto de cooperación en una compañía de automóviles mientras asiste a la escuela. Su supervisora lo está ayudando a aprender sobre el funcionamiento de un motor de combustión interna. Ella le asigna la siguiente tarea, relacionada con la simulación de un nuevo motor que está diseñando. Un gas, comenzando en $P_A = 1.00\text{ atm}$, $V_A = 0.500\text{ L}$ y $T_A = 27.0^\circ\text{C}$, se comprime desde el punto $A$ en el diagrama $PV$ en la figura P7.31 hasta el punto $B$. Esto representa la carrera de compresión en un motor de gasolina de cuatro tiempos. En ese punto, se entregan 132 J de energía al gas a volumen constante, llevando el gas al punto $C$. Esto representa la transformación de la energía potencial en la gasolina en energía interna cuando se enciende la bujía. Su supervisora le dice que la energía interna de un gas es proporcional a la temperatura (como se verá en el capítulo 8). La energía interna del gas en el punto $A$ es 200 J, y ella quiere saber cuál es la temperatura del gas en el punto $C$.
